%!TeX root=../treasuretop.tex
\ArtChapter[I Go to Bristol]{chap7head}

\lettrine[lines=5,image=true,findent=2pt]{chap7initI}{t} was longer than the squire imagined ere we were ready for the sea, and none of our first plans—not even Dr~Livesey's, of keeping me beside him—could be carried out as we intended. The doctor had to go to London for a physician to take charge of his practice; the squire was hard at work at Bristol; and I lived on at the hall under the charge of old Redruth, the gamekeeper, almost a prisoner, but full of sea-dreams and the most charming anticipations of strange islands and adventures. I brooded by the hour together over the map, all the details of which I well remembered. Sitting by the fire in the housekeeper's room, I approached that island in my fancy from every possible direction; I explored every acre of its surface; I climbed a thousand times to that tall hill they call the Spy-glass, and from the top enjoyed the most wonderful and changing prospects. Sometimes the isle was thick with savages, with whom we fought, sometimes full of dangerous animals that hunted us, but in all my fancies nothing occurred to me so strange and tragic as our actual adventures.

So the weeks passed on, till one fine day there came a letter addressed to Dr~Livesey, with this addition, »To be opened, in the case of his absence, by Tom Redruth or young Hawkins.« Obeying this order, we found, or rather I found—for the gamekeeper was a poor hand at reading anything but print—the following important news:

\begin{quotation}
	\begin{flushright}\scshape
Old Anchor Inn, Bristol\\March 1, 17—\\
\end{flushright}
~\\
\nopagebreak
\textsc{Dear Livesey}—As I do not know whether you are at the hall or still in London, I send this in double to both places.

The ship is bought and fitted. She lies at anchor, ready for sea. You never imagined a sweeter schooner—a child might sail her—two hundred tons; name, \textit{Hispaniola}.

I got her through my old friend, Blandly, who has proved himself throughout the most surprising trump. The admirable fellow literally slaved in my interest, and so, I may say, did everyone in Bristol, as soon as they got wind of the port we sailed for—treasure, I mean.
\end{quotation}

»Redruth,« said I, interrupting the letter, »Dr. Livesey will not like that. The squire has been talking, after all.«

»Well, who's a better right?« growled the gamekeeper. »A pretty rum go if squire ain't to talk for Dr~Livesey, I should think.«

At that I gave up all attempts at commentary and read straight on:

\begin{quotation}
Blandly himself found the \textit{Hispaniola}, and by the most admirable management got her for the merest trifle. There is a class of men in Bristol monstrously prejudiced against Blandly. They go the length of declaring that this honest creature would do anything for money, that the \textit{Hispaniola} belonged to him, and that he sold it me absurdly high—the most transparent calumnies. None of them dare, however, to deny the merits of the ship.

So far there was not a hitch. The workpeople, to be sure—riggers and what not—were most annoyingly slow; but time cured that. It was the crew that troubled me.

I wished a round score of men—in case of natives, buccaneers, or the odious French—and I had the worry of the deuce itself to find so much as half a dozen, till the most remarkable stroke of fortune brought me the very man that I required.

I was standing on the dock, when, by the merest accident, I fell in talk with him. I found he was an old sailor, kept a public-house, knew all the seafaring men in Bristol, had lost his health ashore, and wanted a good berth as cook to get to sea again. He had hobbled down there that morning, he said, to get a smell of the salt.

I was monstrously touched—so would you have been—and, out of pure pity, I engaged him on the spot to be ship's cook. Long John Silver, he is called, and has lost a leg; but that I regarded as a recommendation, since he lost it in his country's service, under the immortal Hawke. He has no pension, Livesey. Imagine the abominable age we live in!

Well, sir, I thought I had only found a cook, but it was a crew I had discovered. Between Silver and myself we got together in a few days a company of the toughest old salts imaginable—not pretty to look at, but fellows, by their faces, of the most indomitable spirit. I declare we could fight a frigate.


Long John even got rid of two out of the six or seven I had already engaged. He showed me in a moment that they were just the sort of fresh-water swabs we had to fear in an adventure of importance.

I am in the most magnificent health and spirits, eating like a bull, sleeping like a tree, yet I shall not enjoy a moment till I hear my old tarpaulins tramping round the capstan. Seaward, ho! Hang the treasure! It's the glory of the sea that has turned my head. So now, Livesey, come post; do not lose an hour, if you respect me.

Let young Hawkins go at once to see his mother, with Redruth for a guard; and then both come full speed to Bristol.
\begin{flushright}
\textsc{John Trelawney}
\end{flushright}
%\enlargethispage{2\baselineskip}


\textsc{Postscript}—I did not tell you that Blandly, who, by the way, is to send a consort after us if we don't turn up by the end of August, had found an admirable fellow for sailing master—a stiff man, which I regret, but in all other respects a treasure. Long John Silver unearthed a very competent man for a mate, a man named Arrow. I have a boatswain who pipes, Livesey; so things shall go man-o'-war fashion on board the good ship \textit{Hispaniola}.

I forgot to tell you that Silver is a man of substance; I know of my own knowledge that he has a banker's account, which has never been overdrawn. He leaves his wife to manage the inn; and as she is a woman of colour, a pair of old bachelors like you and I may be excused for guessing that it is the wife, quite as much as the health, that sends him back to roving.
\begin{flushright}
\textsc{J.T.}
\end{flushright}

\textsc{P.P.S.}—Hawkins may stay one night with his mother.\\
\begin{flushright}
\textsc{J.T.}
\end{flushright}

\end{quotation}

You can fancy the excitement into which that letter put me. I was half beside myself with glee; and if ever I despised a man, it was old Tom Redruth, who could do nothing but grumble and lament. Any of the under-gamekeepers would gladly have changed places with him; but such was not the squire's pleasure, and the squire's pleasure was like law among them all. Nobody but old Redruth would have dared so much as even to grumble.

The next morning he and I set out on foot for the »Admiral Benbow«, and there I found my mother in good health and spirits. The captain, who had so long been a cause of so much discomfort, was gone where the wicked cease from troubling. The squire had had everything repaired, and the public rooms and the sign repainted, and had added some furniture—above all a beautiful armchair for mother in the bar. He had found her a boy as an apprentice also so that she should not want help while I was gone.

It was on seeing that boy that I understood, for the first time, my situation. I had thought up to that moment of the adventures before me, not at all of the home that I was leaving; and now, at sight of this clumsy stranger, who was to stay here in my place beside my mother, I had my first attack of tears. I am afraid I led that boy a dog's life, for as he was new to the work, I had a hundred opportunities of setting him right and putting him down, and I was not slow to profit by them.

 

The night passed, and the next day, after dinner, Redruth and I were afoot again and on the road. I said good-bye to Mother and the cove where I had lived since I was born, and the dear old »Admiral Benbow«—since he was repainted, no longer quite so dear. One of my last thoughts was of the captain, who had so often strode along the beach with his cocked hat, his sabre-cut cheek, and his old brass telescope. Next moment we had turned the corner and my home was out of sight.

\begin{figure}[p]
\centering
\includegraphics[width=\linewidth]{chap7mothergoodbyeimproved}
\caption[I said good-bye to Mother]{I said good-bye to Mother and the cove, and the dear old »Admiral Benbow«}
\end{figure}


The mail picked us up about dusk at the »Royal George« on the heath. I was wedged in between Redruth and a stout old gentleman, and in spite of the swift motion and the cold night air, I must have dozed a great deal from the very first, and then slept like a log up hill and down dale through stage after stage, for when I was awakened at last it was by a punch in the ribs, and I opened my eyes to find that we were standing still before a large building in a city street and that the day had already broken a long time.

»Where are we?« I asked.

»Bristol,« said Tom. »Get down.«

Mr~Trelawney had taken up his residence at an inn far down the docks to superintend the work upon the schooner. Thither we had now to walk, and our way, to my great delight, lay along the quays and beside the great multitude of ships of all sizes and rigs and nations. In one, sailors were singing at their work, in another there were men aloft, high over my head, hanging to threads that seemed no thicker than a spider's. Though I had lived by the shore all my life, I seemed never to have been near the sea till then. The smell of tar and salt was something new. I saw the most wonderful figureheads, that had all been far over the ocean. I saw, besides, many old sailors, with rings in their ears, and whiskers curled in ringlets, and tarry pigtails, and their swaggering, clumsy sea-walk; and if I had seen as many kings or archbishops I could not have been more delighted.

And I was going to sea myself, to sea in a schooner, with a piping boatswain and pig-tailed singing seamen, to sea, bound for an unknown island, and to seek for buried treasure!

While I was still in this delightful dream, we came suddenly in front of a large inn and met Squire Trelawney, all dressed out like a sea-officer, in stout blue cloth, coming out of the door with a smile on his face and a capital imitation of a sailor's walk.

»Here you are,« he cried, »and the doctor came last night from London. Bravo! The ship's company complete!«

»Oh, sir,« cried I, »when do we sail?«

»Sail!« says he. »We sail tomorrow!«
